\section{Harmonic heights and connected minima}
\label{sec:heights}

Our goal is to put a height on the spherical cosets inside \(A_Yx\), for
arbitrary \(Y\subseteq S\), whose minimum set is connected.  A stabilizer
preserves that set, so its connectedness will constrain the colors in
which the stabilizer can act.  Lowering a path requires connected choices
for each local descent.  Inside a spherical coset these choices are
described by strict inequalities against the harmonic prescription.
On an isolated null component we must also allow equality; the
null-layer obstruction supplies the nonemptiness needed for this closed
sublevel.  We establish both local connectivity statements before
defining the height.

Only connectedness by finite paths is needed.  Accordingly, every poset
in this section is used through its undirected comparability graph:
two distinct vertices are adjacent when they are comparable.  For coset
vertices, comparability always means both nested types and actual coset
inclusion.

\subsection{Spherical sublevels}

We first record the elementary filtration principle used below.

\begin{lemma}[Connected attaching neighborhoods]
\label{ht:vertex-filtration}
Let a poset be well-ordered by an auxiliary order.  Suppose that, for
each vertex, its earlier neighbors are empty or connected in the graph
on earlier vertices.  Then each connected component contains exactly
one vertex with no earlier neighbor.
\end{lemma}

\begin{proof}
Insert the vertices in the auxiliary order.  A vertex with no earlier
neighbor starts a new component.  Otherwise all its earlier neighbors
belong to one existing component, so insertion enlarges that component
and cannot merge two components.  At a limit stage, any finite path has
already appeared at an earlier stage: its finitely many vertices have
indices bounded below the limit.  Thus components cannot merge at a
limit either.  Every component has a least vertex, which has no earlier
neighbor, proving the assertion.
\end{proof}

Fix a spherical \(T\subseteq S\), a right coset \(A_Tx\), and a cap at an
integer \(q\) on \(\partial T\).  Let \(b^*\) be the harmonic interpolation
of \(m^q(x)\) on \(T\):
\[
 b_i^*=m_i^q(x)\quad(i\notin T),\qquad (Cb^*)_T=0.
\]
Exterior invariance makes \(b^*\) independent of the representative.
Define \(\mathcal R_{<}(T,A_Tx;q)\) to consist of pairs
\[
 v=(J,A_Jz),\qquad J\subsetneq T,\quad z\in A_Tx,
\]
such that every open color \(i\in T\setminus J\) is capped at \(q\) and
has \(m_i^q(z)<b_i^*\).

For a second sublevel, suppose additionally that
\begin{equation}
 \begin{gathered}
 T\subseteq I\subseteq S,\quad I\ne\varnothing\ \text{connected},
 \quad c_I>0,\quad C_Ic_I=0,\\
 c=0\ \text{outside }I,\qquad
 m^q=c\ \text{on }(I\setminus T)\cup\partial I,
 \quad\text{with a cap at }q\text{ there}.
 \end{gathered}
 \label{ht:closed-data}
\end{equation}
Then \(b_T^*=c_T\).  The \emph{special closed sublevel}
\(\mathcal R_{\le}(T,A_Tx;q,I,c)\) has the same definition except that its
open inequalities are \(m_i^q(z)\le c_i\).  In both definitions the type is
proper, and a cap bounds every higher level, not merely level \(q\).

\begin{proposition}[Spherical sublevel connectivity]
\label{ht:sublevels}
The strict sublevel is empty or connected.  Under
Equation~\eqref{ht:closed-data}, the special closed sublevel is nonempty
and connected.  For \(T=\varnothing\), the strict sublevel is empty and
the special closed exterior data cannot occur.  These statements hold
for every integer \(q\).
\end{proposition}

\begin{proof}
For \(T=\varnothing\), there is no proper type.  The closed exterior
data would prescribe \(m^q=c\) and a cap on \(I\cup\partial I\), contrary
to the empty-set case of Proposition~\ref{null:obstruction}.

In the nonempty closed case, Lemma~\ref{cap:access} supplies a
representative capped on \(T\), hence on \(I\).  If the closed sublevel
were empty, every representative capped on \(I\) would satisfy
\(m_j^q>c_j\) for every \(j\in T\): a coordinate with \(m_j^q\le c_j\)
could be opened alone, giving type \(T\setminus\{j\}\).
This contradicts Proposition~\ref{null:obstruction}.  Thus this sublevel
is nonempty.

We prove both connectivity assertions by induction on \(|T|\).
For a vertex \(v=(J,A_Jz)\) in either sublevel, let \(p_v\) be the
harmonic interpolation of \(m^q(z)\) on \(J\), and set
\[
 u_v=(p_v)_T-b_T^*.
\]
The open inequalities and nonnegative interpolation give
\begin{equation}
 u_v\le0,\qquad
 (u_v)_J=-C_J^{-1}C_{J,T\setminus J}(u_v)_{T\setminus J},
 \qquad (Cp_v)_T=C_Tu_v.
 \label{ht:local-difference}
\end{equation}
Order vertices first by increasing \(\sum_{i\in T}(u_v)_i\).
At an equal negative sum put larger types first; at sum zero put
smaller types first.  Refine the remaining ties by any well-order.
Only finitely many sums occur: open counts lie in bounded intervals
with finitely many possible values, the exterior is fixed, and there
are finitely many types.  The resulting order is well founded.
Remaining ties are between incomparable vertices, since distinct
comparable cosets of the same type cannot occur.

Consider first \(u=u_v\ne0\), and write \(V=T\setminus J\).  Put
\[
 N=\{i\in V:(C_Tu)_i\ge0\}.
\]
An upper extension has the form \((K,A_Kz)\), with
\(J\subsetneq K\subsetneq T\), and is still allowed.  The harmonic
change is zero outside \(K\) and on \(K\) equals
\begin{equation}
 (p_{(K,A_Kz)}-p_v)_K=-C_K^{-1}(C_Tu)_K.
 \label{ht:local-upper-change}
\end{equation}
The old rows on \(J\) vanish.  If \(K\setminus J\) avoids \(N\),
the nonnegative inverse with positive diagonal makes the sum strictly
increase.  If \(\varnothing\ne K\setminus J\subseteq N\), the change
is nonpositive; a strict decrease makes the extension earlier, and
equality favors its larger type at the negative sum.
Moreover \(N\ne V\), because otherwise \(C_Tu\ge0\) would imply
\(u=C_T^{-1}C_Tu\ge0\).

Every earlier upper extension therefore contains a new color in \(N\).
Removing its other new colors gives a comparable earlier extension
whose type is contained in \(J\cup N\).  These extensions all connect
through the maximum \((J\cup N,A_{J\cup N}z)\).  Thus the earlier upper
part is empty or connected.

An earlier lower neighbor of \(v\) must strictly decrease the sum.
Write it as \((R,A_Rw)\), where \(R\subsetneq J\), \(w\in A_Jz\).
The difference of its harmonic vector and \(p_v\) is zero outside \(J\),
has prescribed values
\[
 m_i^q(w)-(p_v)_i\qquad(i\in J\setminus R),
\]
and is harmonically interpolated on \(R\).
If all these values were nonnegative, the sum could not decrease.
Hence the set of newly opened colors satisfying
\(m_i^q(w)<(p_v)_i\) is nonempty.  Close the other newly opened colors.
This produces a comparable earlier lower vertex in the strict sublevel
\(\mathcal R_{<}(J,A_Jz;q)\), with harmonic bound \(p_v\).

The boundary of \(J\) is capped: its colors lie in \(V\) or
\(\partial T\).  Every vertex of this inner strict sublevel is conversely
an allowed earlier lower vertex in the original sublevel, because
\(p_v\le b^*\) and its harmonic difference from \(p_v\) has strictly
negative sum.  Induction makes the inner sublevel empty or connected.
Since each earlier lower vertex is comparable with a vertex in it,
the whole earlier lower part is empty or connected as well.

If \(u_v=0\), we are in the special closed case.  An upper extension
keeps the same harmonic vector and has larger type, so is later.
Every allowed lower neighbor is earlier, and these neighbors form
\(\mathcal R_{\le}(J,A_Jz;q,I,c)\).  Its exterior data are supplied by
Equation~\eqref{ht:closed-data} together with the equalities and caps on
\(V\).  The empty \(J\) case is impossible by
Proposition~\ref{null:obstruction}; otherwise induction makes this lower
part nonempty and connected.

If both earlier parts are nonempty, every earlier lower vertex is below
every earlier upper vertex, so their union is connected.
Lemma~\ref{ht:vertex-filtration} now applies.
A vertex with no earlier neighbor must have \(u_v\ne0\), \(N=\varnothing\),
and empty inner strict sublevel.  Therefore
\[
 -(Cp_v)_i>0\quad(i\in V),\qquad (Cp_v)_J=0.
\]
At every representative \(w\in A_Jz\) capped on \(T\) one also has
\(m_J^q(w)\ge(p_v)_J\); otherwise opening a failing coordinate alone
would give a vertex in the inner strict sublevel.
Together with the caps on \(V\), these are exactly the candidate
conditions of Proposition~\ref{two:uniqueness}.  There is at most one
such vertex in the fixed spherical residue.
Every component has one by Lemma~\ref{ht:vertex-filtration}, so a
nonempty sublevel is connected.
\end{proof}

\subsection{The global relative height}

Fix \(Y\subseteq S\) and \(x\in A\), without a sphericity assumption on
\(Y\).  Let \(\mathcal P(Y,x)\) be the poset of
\[
 v=(T,A_Tz),\qquad T\subseteq Y\ \text{spherical},\quad z\in A_Yx.
\]
Its graph is connected.  Indeed singleton cosets \((\varnothing,\{z\})\)
are joined to their left-generator translates through the corresponding
rank-one cosets, and every vertex contains a singleton vertex.

Choose \(d_*\le0\) at or below every occupied level of the exterior data
\(m^k(x)|_{\widehat S\setminus Y}\).
These data are fixed throughout \(A_Yx\).
For \(v=(T,A_Tz)\), let \(p^k(v)\) be the harmonic interpolation of
\(m^k(z)\) on \(T\).  It is independent of \(z\) in its coset.

For each \(k\ge d_*\), assign a sign \(\eta_i^k(v)\) as follows.
It is zero if \(p_i^k(v)=0\), and is \(-1\) at positive entries, except
that it is \(+1\) on a whole positive-support component \(I\) when
\begin{equation}
 I\subseteq S,\qquad (Cp^k(v))_I=0,\qquad
 p^d(v)|_{I\cup\partial I}=0\quad\text{for every }d>k.
 \label{ht:null-sign}
\end{equation}
Here support components use adjacency in the full diagram
\(\widehat S\).  Fix an ordering of those coordinates and form the block
\begin{equation}
 B_k(v)=
 \left(
  \sum_{i\in\widehat S}p_i^k(v),\
  (p_i^k(v))_{i\in\widehat S},\
  (\mathbf1_{\{i\in T\}}\eta_i^k(v))_{i\in\widehat S}
 \right).
 \label{ht:block}
\end{equation}
Blocks are ordered lexicographically, with smaller earlier.
The \emph{height} of \(v\) compares these finite-support blocks from the
highest level downward, stopping at \(d_*\).
Distinct vertices may have equal height.

The final coordinates of each block break numerical ties.  When the
recorded interpolated vectors are unchanged, adding a positive color to
the type changes its coordinate from zero to \(\eta_i^k\): sign \(-1\)
favors closing that color, while sign \(+1\) favors opening it.
The exceptional \(+1\) case in Equation~\eqref{ht:null-sign} is an
isolated highest positive null component.  An entire such component is
nonspherical, so it cannot all be added to a spherical type.  Openings
whose counts are at most the prescribed null vector are instead
controlled by the special closed sublevel of
Proposition~\ref{ht:sublevels}.

\begin{lemma}[Well-founded heights]
\label{ht:well-order}
The possible heights are well ordered.  Their least block is the zero
block.
\end{lemma}

\begin{proof}
All entries of \(p^k\) are nonnegative by Lemma~\ref{cat:cartan}.
Bounding a block's first
coordinate bounds every exterior count for its type.  Those counts
have finitely many possible values in a bounded interval by
Proposition~\ref{cat:family}.  There are
finitely many types, their harmonic interpolations are fixed linear
maps, and there are finitely many sign masks.
Thus the block alphabet has finite initial segments.  Its least
element is zero.

An infinite strictly descending sequence of heights would have no
nonzero block above the largest occupied level of its first term.
It would therefore involve only finitely many levels \(k\ge d_*\).
At the highest level the block can decrease only finitely often, so
eventually stays fixed.  The same argument then applies to the next
level, and successively to all remaining levels, contradicting strict
descent.  Equivalently, if \(r(B)\) is the finite rank of a block in its
alphabet, the heights are ordered by the finite Cantor sums
\(\sum_{k\ge d_*}\omega^{k-d_*}r(B_k)\), with decreasing exponents.
\end{proof}

A neighbor is called \emph{descending} when it has strictly smaller
height.  We analyze upper and lower neighbors separately.

\subsection{Active upper colors}

Fix \(v=(T,A_Tz)\).  For \(i\in Y\setminus T\), let \(\ell_i\) be the
highest recorded level with a positive entry on \(i\) or one of its
neighbors, if such a level exists.  Call \(i\) \emph{active} if
\begin{equation}
 (Cp^{\ell_i}(v))_i>0
 \quad\text{or}\quad
 (Cp^{\ell_i}(v))_i=0,\quad \eta_i^{\ell_i}(v)=-1.
 \label{ht:active}
\end{equation}
An active color has a positive entry at that level: otherwise incidence
would make its row strictly negative.

\begin{lemma}[Upper neighbors]
\label{ht:upper}
The union of \(T\) with all active colors is spherical.
The descending upper neighbors of \(v\) are empty or connected using
only descending upper neighbors.
\end{lemma}

\begin{proof}
An upper neighbor has the form \(v'=(U,A_Uz)\), where \(T\subsetneq U\)
is spherical.  Compare at the highest incidence level of its new colors.
Above that level, the new colors and their neighbors have zero entries,
so interpolation changes neither vectors nor signs.
At the comparison level \(k\), the change is zero outside \(U\) and is
\begin{equation}
 (p^k(v')-p^k(v))_U=-C_U^{-1}(Cp^k(v))_U.
 \label{ht:global-upper-change}
\end{equation}
Rows on \(T\) vanish.
If no new color is active, all the new rows at this level are
nonpositive.  A negative row strictly increases the sum.  If all
vanish, the vector and its intrinsic signs stay fixed, and the incident
positive new coordinates have sign \(+1\).  Closing them increases
the sign mask.  If the new colors have no recorded incidence at all,
the heights tie.
Consequently every descending upper neighbor adds an active color.

Conversely, suppose only active colors are added.
At their highest incidence level all new rows are nonnegative.
A positive row strictly decreases the sum; if all vanish, the vector
is fixed and the positive new coordinates have sign \(-1\), so the
sign mask decreases.  Thus such an extension is descending whenever
its type is spherical.

Let \(G\) be a connected component of \(T\) together with all active
colors, and suppose \(G\) contains an active color.
Choose the highest recorded level \(k\) incident to \(G\).
All full rows on \(G\) are nonnegative: closed rows vanish, while an
active color either has its highest incidence at \(k\), or has zero
data on itself and its neighbors there.  Hence
\begin{equation}
 C_Gp_G^k=(Cp^k)_G+b_{G,\widehat S\setminus G}p_{\widehat S\setminus G}^k
 \ge0.
 \label{ht:active-cartan}
\end{equation}
The vector \(p_G^k\) is nonzero, since otherwise incidence would give a
negative full row.  It is strictly positive by connectedness: a zero
coordinate adjacent to a positive one would again have a negative row.

If \(G\) were nonspherical, Lemma~\ref{cat:cartan} would give
\(C_Gp_G^k=0\).  Both nonnegative terms on the right of
Equation~\eqref{ht:active-cartan} would vanish.
Thus \(G\) is the entire positive-support component, and higher levels
vanish on \(G\cup\partial G\) by the choice of \(k\).
Equation~\eqref{ht:null-sign} assigns sign \(+1\) to \(G\).
Every active color in it has highest incidence \(k\) and zero full row,
contradicting Equation~\eqref{ht:active}.  Therefore \(G\) is spherical.
Components with no active color lie in the old spherical type \(T\).
The whole union is consequently spherical.

Every descending upper neighbor can now be joined to the one obtained
by removing all its nonactive new colors.  The latter is descending
and still contains a new color.  All these reduced extensions connect
through the extension by every active color, which is spherical by the
preceding argument.  If there are no active colors, there is no
descending upper neighbor.
\end{proof}

\subsection{Eligible lower colors and null blocks}

A connected diagram component of \(T\) is \emph{visible} if its harmonic
entries are positive at some recorded level.
Whenever they are nonzero at a level, all its entries are positive,
by positivity of the inverse Cartan block.
It therefore has a common highest such level \(q\).
Its boundary data are capped there: at every higher level its zero
harmonic entries force all neighboring nonnegative entries to vanish.
Its entries also have a common sign at level \(q\).

For a lower coset, call a newly opened color \(i\) \emph{eligible} when
it belongs to a visible component with highest level \(q\), its actual
data are capped at \(q\), and
\begin{equation}
 m_i^q<p_i^q(v)
 \quad\text{or}\quad
 m_i^q=p_i^q(v),\quad \eta_i^q(v)=+1.
 \label{ht:eligible}
\end{equation}
Nonvisible colors are not eligible.

\begin{lemma}[Lower neighbors]
\label{ht:lower}
Every descending lower neighbor opens an eligible color.
Opening a nonempty set of eligible colors and no others gives a
descending neighbor.  The descending lower neighbors of \(v\) are
empty or connected using only descending lower neighbors.
\end{lemma}

\begin{proof}
Write a lower neighbor as \((R,A_Rw)\), with
\(R\subsetneq T\) and \(w\in A_Tz\).  Its vector difference from
\(p^k(v)\) is zero outside \(T\), has actual-minus-prescribed values on
the newly opened colors \(T\setminus R\), and is harmonically
interpolated on \(R\).
Consider the highest recorded level at which an opened coordinate has
either prescribed or actual positive data.
Above it both vectors and sign masks agree.

Suppose no opened color is eligible.
At the comparison level the actual-minus-prescribed values are all
nonnegative.  Indeed any positive prescribed coordinate is at its
component's highest visible level, and its actual cap holds, because
there is no higher incidence among the opened coordinates.
Failure of eligibility then gives the nonnegative difference.
A strict difference increases the sum after interpolation.
If all differences vanish, the full vectors and intrinsic signs agree,
and every positive opened coordinate has sign \(-1\).  Its mask changes
from \(-1\) to zero, increasing the block.
Without a comparison level there is a tie.
This proves that a descending lower neighbor must open an eligible
color.

If all opened colors are eligible, the reverse comparison applies.
At the highest visible level among them their differences are
nonpositive, and the other opened coordinates have zero difference.
A strict difference decreases the sum.  With equality, some positive
opened coordinate has sign \(+1\), and every nonzero mask change is
from \(+1\) to zero.  Thus the neighbor is descending.

From any descending lower neighbor, close its ineligible opened
colors.  This gives a comparable descending lower neighbor opening
only eligible colors.  It remains to connect all such choices.

First take a visible component with sign \(-1\) at its highest level
\(q\).  The proper eligible choices in that component are exactly its
strict spherical sublevel, with the fixed exterior cap at \(q\).
This factor is empty or connected by Proposition~\ref{ht:sublevels}.

For sign \(+1\), components must be grouped more carefully.
Let \(I\) be their whole positive-support component at level \(q\);
group together all the diagram components of \(T\) contained in \(I\).
Their union is precisely \(T\cap I\).
No other closed color is in \(\partial I\): a closed zero harmonic
entry cannot neighbor positive support.
The sign rule says that \(I\subseteq S\), that its positive vector
\(c_I=p_I^q(v)\) satisfies \(C_Ic_I=0\), and that all higher incident
data vanish.
Since \((I\setminus T)\cup\partial I\) is outside \(T\), its harmonic
data are the fixed actual data.  They therefore satisfy
Equation~\eqref{ht:closed-data} for the spherical block \(T\cap I\).
Its proper eligible choices form the special closed sublevel for
this entire block, which is nonempty and connected by
Proposition~\ref{ht:sublevels}.  No separate closed-sublevel assertion
for its possibly disconnected diagram components is being used.

All these blocks are unions of components of \(T\); nonvisible
components stay closed.  The Artin presentation and standard
parabolic embedding give the direct product of the block groups.
Thus a subcoset in \(A_Tz\) is specified by its block subcosets acting
from the common representative \(z\).
Operations in another block change no counts in a given block.
Consequently the all-eligible lower choices are exactly the product
of the proper-choice factors just described, each with its top
(the choice to open nothing) adjoined, with the all-top tuple deleted.

This product is empty or connected.
Every tuple can be joined upward to one with a single non-top
coordinate.  Such tuples connect within a fixed coordinate because
its proper-choice factor is connected.  Tuples supported in two
different coordinates connect through the tuple carrying both
choices.  Empty factors merely force their coordinates to stay top.
Hence all-eligible neighbors are empty or connected; adjoining the
other descending neighbors through their comparable all-eligible
choices proves the result.
\end{proof}

\begin{corollary}[Descending neighborhoods]
\label{ht:descending}
The descending neighbors of each vertex are empty or connected in its
strict height sublevel.
\end{corollary}

\begin{proof}
Lemmas~\ref{ht:upper} and~\ref{ht:lower} handle the two parts.
If both are nonempty, each lower neighbor is below each upper neighbor,
so they are joined by edges.
\end{proof}

\subsection{Ties and minimum-height paths}

For \(v=(T,A_Tz)\), define its invisible and visible closed sets by
\begin{equation}
 \begin{split}
 K(v)&=\{i\in Y:
  p^k(v)|_{\{i\}\cup\partial\{i\}}=0
       \text{ for every }k\ge d_*\},\\
 T_+(v)&=T\setminus K(v).
 \end{split}
 \label{ht:invisible}
\end{equation}
These definitions do not assert absence of internal generators below
the recorded cutoff.

\begin{lemma}[Compatibility across a tie]
\label{ht:ties}
If \(v<v'\) have equal height, then \(K(v)=K(v')\) and
\(T_+(v)=T_+(v')\).
There is no adjacency between these common visible and invisible sets.
Moreover either both descending neighborhoods are empty, or they are
nonempty and belong to the same connected component of their common
strict height sublevel.
\end{lemma}

\begin{proof}
Write \(v=(T,A_Tz)\) and \(v'=(U,A_Uz)\), with \(T\subsetneq U\).
All recorded vectors agree.  Equality of each sign-mask coordinate
forces every new closed color in \(U\setminus T\) to have zero entry
at every recorded level: at a positive entry its mask would change
from zero to \(+1\) or \(-1\).
Harmonicity on \(U\) then forces all its neighboring entries to vanish
too.  Hence \(U\setminus T\subseteq K(v)\), the invisible sets agree,
and
\[
 T\setminus K(v)=U\setminus K(v).
\]
A visible closed color is positive at some recorded level; it therefore
cannot neighbor an invisible color.  This proves the first assertions.

Active colors are common, because they depend only on the recorded
vectors and their signs; no new closed color is active.
If there is an active color, its upper extension from the larger
coset is descending and comparable with both vertices.
Corollary~\ref{ht:descending} then places both entire descending
neighborhoods in the same strict-sublevel component.

Suppose instead that there are no active colors.
Both upper descending parts are empty, so only lower neighbors matter.
Their visible components, highest levels, signs, and grouped
eligibility blocks agree.
If the larger coset has a descending lower neighbor, it has an
eligible opening after a move \(g\in A_U\).
The nonadjacency already proved gives
\[
 A_U=A_{T_+}\times A_{U\cap K}.
\]
Write \(g=g_+g_0\) in these factors.
Omitting \(g_0\) cannot change any actual count at a visible color:
it commutes with \(g_+\) and operates only in invisible colors.
Since \(g_+\in A_T\), the same eligible color can be opened in the
smaller coset, using the representative \(g_+z\).
Opening that single color gives a descending lower neighbor of the
smaller vertex by Lemma~\ref{ht:lower}.
Conversely every descending lower neighbor of the smaller vertex
is also one of the larger, by transitivity of coset inclusion and
equality of their heights.
Thus nonemptiness agrees and, when nonempty, the two neighborhoods
share a lower vertex.
Corollary~\ref{ht:descending} finishes the proof.
\end{proof}

\begin{proposition}[Connected minima]
\label{ht:minima}
The minimum-height vertices in \(\mathcal P(Y,x)\) form a nonempty
connected graph.  Their sets \(K\) and \(T_+\) are constant along this
graph.
\end{proposition}

\begin{figure}[htbp]
\centering
\begin{tikzpicture}[x=1cm,y=1cm,
  vertex/.style={circle,fill=black,inner sep=1.8pt},
  lower/.style={circle,fill=linkcolor,inner sep=1.8pt},
  every node/.style={font=\small}]
  \draw[densely dashed,gray] (-.7,1.35)--(7.6,1.35);
  \node[anchor=east] at (-.8,1.35) {height $H$};
  \node[anchor=east,align=right] at (-.8,-.15) {strict sublevel\\of $H$};
  \coordinate (a) at (0,-.1);
  \coordinate (b) at (6.6,-.1);
  \coordinate (v0) at (1.2,1.35);
  \coordinate (v1) at (3.3,1.35);
  \coordinate (v2) at (5.4,1.35);
  \draw[thick] (a)--(v0)--(v1)--(v2)--(b);
  \node[vertex,label=above:$v_0$] at (v0) {};
  \node[vertex,label=above:$v_1$] at (v1) {};
  \node[vertex,label=above:$v_2$] at (v2) {};
  \node at (3.3,2.02) {maximal tied segment};
  \coordinate (u0) at (1.4,-.8);
  \coordinate (u1) at (3.3,-.45);
  \coordinate (u2) at (5.2,-.8);
  \draw[thick,linkcolor] (a)--(u0)--(u1)--(u2)--(b);
  \foreach \p in {u0,u1,u2} {\node[lower] at (\p) {};}
  \node[vertex,label=left:$a$] at (a) {};
  \node[vertex,label=right:$b$] at (b) {};
  \node[text=linkcolor] at (3.3,-1.2) {replacement path};
\end{tikzpicture}

\caption{Lowering a maximal tied segment, as in
Proposition~\ref{ht:minima}.  Its two descending attachments $a,b$
belong to the same component of the strict height sublevel, so a
finite path there replaces the segment.  The drawing represents
comparability paths and a connectivity argument.}
\label{fig:tied-path}
\end{figure}

\begin{proof}
A minimum exists by Lemma~\ref{ht:well-order}.
Take a finite path between two minimum-height vertices in the
connected full graph.
If its greatest height is above the minimum, consider a maximal
contiguous segment at that height.
Consecutive vertices of the segment are comparable and tied.
The segment has descending attachments at its two ends.
Lemma~\ref{ht:ties} and Corollary~\ref{ht:descending} put these attachments
in one connected component of the strict sublevel.
Replace the segment by a finite path in that component.

Doing this for all segments at the current greatest height removes
that height entirely and introduces no vertex at or above it.
If the resulting path is still not a minimum path, repeat.
An infinite repetition would give a strictly descending sequence of
greatest heights, contrary to Lemma~\ref{ht:well-order}.
The process therefore ends with a path entirely among minima.
The constancy of \(K,T_+\) along its edges is Lemma~\ref{ht:ties}.
\end{proof}

\subsection{The stabilizer containment}

For later use define the simultaneous graded-homotopy stabilizer
\begin{equation}
 H(Y,x)=\{a\in A_Y:
 E_L(ax)\simeq E_L(x)\text{ for every simple label }L\}.
 \label{ht:stabilizer}
\end{equation}
This is a subgroup because the \(B_a\) are equivalences.
Choose a minimum \(v=(T,A_Tz)\), and abbreviate its invisible and visible
sets by \(K,T_+\).

\begin{proposition}[Containment from a minimum]
\label{ht:containment}
Put \(D=T_+\cup K\subseteq Y\).
The set \(T_+\) is spherical and nonadjacent to \(K\), and
\begin{equation}
 x^{-1}H(Y,x)x\subseteq z^{-1}A_Dz.
 \label{ht:containment-equation}
\end{equation}
Equivalently, with \(b=zx^{-1}\in A_Y\),
\(H(Y,x)\subseteq b^{-1}A_Db\).
\end{proposition}

\begin{proof}
The set \(T_+\) is contained in the spherical type \(T\).
Every visible closed color is positive at some recorded level, so
Equation~\eqref{ht:invisible} excludes adjacency to \(K\).

Along any minimum path starting at \(v\), the sets \(T_+,K,D\) are
constant by Proposition~\ref{ht:minima}, and every type is contained
in \(D\).
All cosets on that path lie in \(A_Dz\).
Indeed this holds initially because \(T\subseteq D\).
A downward inclusion preserves it; for an upward inclusion, the larger
coset has type contained in \(D\) and contains a point of \(A_Dz\), so
the whole larger coset also lies in \(A_Dz\).

Let \(a\in H(Y,x)\) and set \(r=x^{-1}ax\).
Right multiplication by \(r\) is an automorphism of the coset poset,
preserving types.  At a representative \(cx\), where \(c\in A_Y\),
\[
 E_L(cxr)=E_L(cax)\simeq B_cE_L(ax)
 \simeq B_cE_L(x)=E_L(cx).
\]
It therefore preserves all layer counts and heights.
The translated vertex \((T,A_Tzr)\) is a minimum, so
Proposition~\ref{ht:minima} joins it to \(v\) by a minimum path.
The coset containment along that path gives \(zr\in A_Dz\), hence
\(r\in z^{-1}A_Dz\), proving Equation~\eqref{ht:containment-equation}.
The equivalent formulation follows from \(z=bx\).
\end{proof}
